\begin{theorem}[Different exponent of $\mathbb{Q}_2(\sqrt2)/\mathbb{Q}_2$]
Let $L = \mathbb{Q}_2(\sqrt2) = \mathbb{Q}_2[X]/(X^2-2)$. Then the different
exponent of $L/\mathbb{Q}_2$ equals $3$, i.e.
\[
  \delta(L/\mathbb{Q}_2)
  \;=\; v_{\mathfrak{m}_L}\!\bigl(\mathfrak{d}_{L/\mathbb{Q}_2}\bigr)
  \;=\; 3 .
\]
\end{theorem}

\begin{proof}
Write $K = \mathbb{Q}_2$, $\mathcal{O}_K = \mathbb{Z}_2$, and let
$\mathcal{O}_L$ be the ring of integers of $L$. Let $\theta \in \mathcal{O}_L$
be the image of the canonical root of $X^2 - 2$, so that
\begin{equation}\label{eq:theta-sq}
  \theta^2 = 2,\qquad \theta \in \mathfrak{m}_L .
\end{equation}
Recall from the analysis of the Eisenstein extension $\mathbb{Q}_p(p^{1/e})$
(applied here with $p = e = 2$) that
\begin{equation}\label{eq:ef}
  e(L/K) = 2,\qquad f(L/K) = 1,\qquad \mathfrak{m}_K\mathcal{O}_L=\mathfrak{m}_L^{\,2}.
\end{equation}

\medskip
\noindent\emph{Step 1: $\theta$ is a uniformizer of $\mathcal{O}_L$.}
By \eqref{eq:theta-sq} and \eqref{eq:ef} the element $\theta^2 = 2$ generates
$\mathfrak{m}_K\mathcal{O}_L = \mathfrak{m}_L^{2}$, so
$(\theta)^2 = (\theta^2) = \mathfrak{m}_L^{2}$. Since $\mathcal{O}_L$ is a
discrete valuation ring, the map $n \mapsto \mathfrak{m}_L^{\,n}$ is strictly
decreasing on ideals; hence $(\theta)^2 = \mathfrak{m}_L^2$ forces
$(\theta) = \mathfrak{m}_L$. Thus $\theta$ is a uniformizer, i.e. irreducible in
$\mathcal{O}_L$.

\medskip
\noindent\emph{Step 2: $\mathcal{O}_L = \mathcal{O}_K[\theta]$.}
Because $f(L/K) = 1$, the residue extension $\lambda/\kappa$ is trivial, so any
residue class is generated by (the image of) $\theta$; equivalently the residue
field is generated over $\kappa$ by $\bar\theta$. Combined with Step~1
($\theta$ a uniformizer), the two-generator monogenicity criterion (taking both
generators equal to $\theta$) gives
\[
  \mathcal{O}_K[\theta] = \mathcal{O}_L .
\]

\medskip
\noindent\emph{Step 3: the minimal polynomial.}
The polynomial $g = X^2 - 2 \in \mathcal{O}_K[X]$ is monic of degree $2$ and,
by \eqref{eq:theta-sq}, annihilates $\theta$. Since $[L:K] = 2$ and
$\theta \notin K$, the minimal polynomial of $\theta$ over $K$ has degree $2$;
as $\mathcal{O}_K$ is integrally closed, the minimal polynomial of $\theta$ over
$\mathcal{O}_K$ is obtained by clearing denominators and equals the monic
degree-$2$ polynomial annihilating $\theta$, namely
\[
  \operatorname{minpoly}_{\mathcal{O}_K}(\theta) = g = X^2 - 2 .
\]

\medskip
\noindent\emph{Step 4: the different is generated by $g'(\theta)$.}
By Step~2 the extension is monogenic, so the conductor of $\mathcal{O}_K[\theta]$
in $\mathcal{O}_L$ is the unit ideal. The single-generator formula for the
different then yields
\[
  \mathfrak{d}_{L/K}
  = \bigl(g'(\theta)\bigr)
  = \bigl(2\theta\bigr).
\]

\medskip
\noindent\emph{Step 5: the collapse $2\theta = \theta^3$.}
Using \eqref{eq:theta-sq},
\[
  2\theta = \theta^2\cdot\theta = \theta^3 .
\]
Therefore, by Step~1,
\[
  \mathfrak{d}_{L/K} = (2\theta) = (\theta^3) = (\theta)^3 = \mathfrak{m}_L^{\,3}.
\]

\medskip
\noindent\emph{Step 6: reading off the exponent.}
In the discrete valuation ring $\mathcal{O}_L$ the maximal ideal $\mathfrak{m}_L$
is nonzero and not a unit, so the $\mathfrak{m}_L$-adic multiplicity of
$\mathfrak{m}_L^{\,3}$ is exactly $3$. Hence
\[
  \delta(L/\mathbb{Q}_2)
  = v_{\mathfrak{m}_L}\!\bigl(\mathfrak{d}_{L/K}\bigr)
  = v_{\mathfrak{m}_L}\!\bigl(\mathfrak{m}_L^{\,3}\bigr)
  = 3 . \qedhere
\]
\end{proof}